\chapter{RLHF: Reward Models, Value Functions, and PPO}\label{ch:rlhf}

Chapter~\ref{ch:pg} derived the policy gradient and the tools that tame its variance:
baselines, advantages, importance ratios and clipping. This chapter puts those tools to work in
the system that made chat models possible: \emph{reinforcement learning from human feedback}
(RLHF) with Proximal Policy Optimization (PPO). Chapter~\ref{ch:grpo} met RLHF briefly as the
ancestor of RLVR; here we build all of it. That means collecting preference data, training a reward model
on it, adding a value model (the critic), shaping per-token rewards with a KL penalty,
estimating advantages with GAE, and running the full PPO loop over four interacting networks.

For a systems engineer, PPO-RLHF is the most demanding training workload in this book. One
iteration generates text autoregressively, which is memory-bound and runs at small batch. It then
runs prefill-style forward passes through four models, and finally performs several epochs of
backward passes through two of them. All four models are as large as the policy or close to it.
\Cref{sec:rlhf:system} works out where the memory and time go and how real frameworks place the
models. The lab builds the complete system at toy scale and shows, on a laptop CPU, the
two things every practitioner must see at least once: PPO raising a learned reward, and the policy
\emph{over-optimising} it, so that the proxy reward keeps rising while true quality falls.

\begin{figure}[t]
\centering
\resizebox{\linewidth}{!}{%
\begin{tikzpicture}[
  st/.style={draw, rounded corners=3pt, align=center, font=\small, minimum height=11mm,
    text width=25mm, line width=0.6pt},
  num/.style={circle, fill=sysblue, text=white, font=\scriptsize\bfseries, inner sep=1.5pt},
  >={Stealth[length=2.2mm]}]
\node[st, fill=fdata!15] (s1) at (0,0) {Preference data\\{\scriptsize pairs, rankings}};
\node[st, fill=napkin!18] (s2) at (3.3,0) {Reward model\\{\scriptsize BT loss, $r_\phi$}};
\node[st, fill=fcompute!15] (s3) at (6.6,0) {Per-token rewards\\{\scriptsize score $-\beta\,$KL}};
\node[st, fill=fcompute!15] (s4) at (9.9,0) {Value model + GAE\\{\scriptsize $V_\psi$, $\hat A_t$}};
\node[st, fill=fcompute!25] (s5) at (13.2,0) {PPO update\\{\scriptsize clip, value loss}};
\node[st, fill=fcomm!15] (s6) at (13.2,-2.1) {Diagnostics\\{\scriptsize KL, EV, over-opt.}};
\node[st, fill=sysbluebg] (s7) at (9.9,-2.1) {Four-model system\\{\scriptsize memory, placement}};
\node[st, fill=keygreen!12] (s8) at (6.6,-2.1) {RLAIF, agents,\\case studies};
\foreach \n/\k in {s1/1,s2/2,s3/3,s4/4,s5/5,s6/6,s7/7,s8/8}
  \node[num] at (\n.north west) {\k};
\draw[->] (s1)--(s2); \draw[->] (s2)--(s3); \draw[->] (s3)--(s4); \draw[->] (s4)--(s5);
\draw[->] (s5)--(s6); \draw[->] (s6)--(s7); \draw[->] (s7)--(s8);
\draw[->, dashed, fcompute] (s5.north) -- ++(0,0.5) -| node[pos=0.25, above, font=\scriptsize]
  {new rollouts from the updated policy} (s3.north);
\end{tikzpicture}}
\caption{Roadmap of the chapter. Stages 1--2 turn human judgements into a differentiable scorer;
stages 3--5 form the PPO loop, which runs many times (dashed arrow); stages 6--8 are about
making it work and making it fit on hardware.}
\label{fig:rlhf:roadmap}
\end{figure}

\section{Why a third training stage?}\label{sec:rlhf:why}

Pretraining (\cref{ch:pretrain}) produces a model of text, and supervised fine-tuning
(\cref{ch:instruct}) imitates demonstrations of the behaviour we want. Imitation has two limits
that a third stage removes.

\begin{enumerate}
  \item \textbf{Demonstrations are expensive and capped by the demonstrator.} Writing an excellent
        answer is far harder than recognising one. This \emph{generator--verifier gap} is the
        economic argument for RLHF: annotators compare two model outputs in seconds, but they
        would need minutes to write a better one.
  \item \textbf{SFT has no negative signal.} Maximum likelihood pushes probability \emph{up} on the
        demonstrations and has no direct way to say ``never do \emph{this}''. A comparison
        ``$A$ is better than $B$'' carries information about both responses.
\end{enumerate}

The RLHF stage therefore solves the KL-regularised problem
\begin{equation}
  \max_{\theta}\; \E_{x\sim\mathcal{D},\,y\sim\pitheta(\cdot\mid x)}
  \big[r_\phi(x,y)\big] \;-\; \beta\,\E_{x}\,
  \KL\!\big(\pitheta(\cdot\mid x)\,\big\|\,\piref(\cdot\mid x)\big),
  \label{eq:rlhf:objective}
\end{equation}
where $r_\phi$ is a \emph{learned} reward model, $\piref$ is the SFT model, and $\beta$ sets
how far the policy may wander from it. Each term in \cref{eq:rlhf:objective} is needed. Without
the reward, nothing changes. Without the KL term, the policy finds inputs on which the learned
reward is wrong and exploits them, a failure we will reproduce in \cref{sec:rlhf:overopt}.
Chapter~\ref{ch:dpo} shows that the same objective admits a closed-form optimum, which leads to
DPO. Here we optimise it the original way, with on-policy RL.

\begin{keyidea}[The three-stage recipe]
InstructGPT (Ouyang et al., 2022) set the template that most later chat models follow: \textbf{(1)}
SFT on demonstrations, \textbf{(2)} a reward model trained on human comparisons of SFT outputs,
\textbf{(3)} PPO against that reward model with a KL penalty to the SFT model. Later work varies
every stage. DPO (\cref{ch:dpo}) skips the explicit reward model, RLVR (\cref{ch:grpo}) swaps it
for a program, and GRPO drops the critic. All of them, however, keep the
generate--score--update loop built in this chapter.
\end{keyidea}

\section{From classical RL to RLHF: the actor--critic view}\label{sec:rlhf:actorcritic}


\Cref{ch:pg} introduced the mapping between RL and language modelling: the state is the prompt
plus the tokens generated so far, an action is the next token, the environment's transition is
deterministic concatenation, and an episode is one response. PPO adds one component that GRPO
does not have, a learned \emph{critic}. RL algorithms fall into three families:

\begin{table}[h]
\centering\small
\begin{tabularx}{\linewidth}{@{}lXXl@{}}
\toprule
\textbf{Family} & \textbf{What is learned} & \textbf{How actions are chosen} & \textbf{LLM examples}\\
\midrule
Value-based & $Q(s,a)$ or $V(s)$ & $\argmax_a Q(s,a)$ & rare for LLMs (huge action space)\\
Policy-based & $\pitheta(a\mid s)$ & sample from the policy & REINFORCE, GRPO, RLOO\\
Actor--critic & policy \emph{and} $V_\psi(s)$ & sample; critic supplies the baseline & PPO\\
\bottomrule
\end{tabularx}
\caption{Where PPO sits. DPO is policy-based but is not a policy-gradient method at all: it is a
supervised loss on preference pairs.}
\label{tab:rlhf:families}
\end{table}

In PPO the \textbf{actor} is the policy LLM and the \textbf{critic} is a second network
$V_\psi(s_t)$ that predicts, for every prefix, the return the actor will eventually collect.
Its job is to turn one noisy end-of-sequence reward into a per-token learning signal, the
advantage $\hat A_t$: did choosing token $t$ make the outcome better or worse than expected?
\Cref{tab:rlhf:mapping} lists what changes when classical PPO is moved onto text.

\begin{table}[h]
\centering\small
\begin{tabularx}{\linewidth}{@{}lXX@{}}
\toprule
& \textbf{Classical PPO (e.g.\ robotics, games)} & \textbf{PPO in RLHF}\\
\midrule
Environment & physics simulator or game engine & string concatenation (deterministic)\\
Episode & hundreds to millions of steps & one response, tens to thousands of tokens\\
Reward & dense, from the environment & one scalar from $r_\phi$ at the last token, plus a per-token KL penalty\\
Discount $\gamma$ & $0.99$ & usually $1$ (episodes are short and finite)\\
Action space & small, or a low-dimensional continuous vector & vocabulary, $|V|\approx 10^5$\\
Critic & small MLP & a full transformer, as large as the policy or nearly\\
Reference policy & none & SFT model $\piref$ (frozen), anchors the KL term\\
Cost of an episode & cheap to simulate & autoregressive decoding of a large model\\
\bottomrule
\end{tabularx}
\caption{The same algorithm, a very different workload. The last two rows explain most of the
systems difficulty of PPO-RLHF.}
\label{tab:rlhf:mapping}
\end{table}

\section{Collecting preference data}\label{sec:rlhf:data}

The reward model can only be as good as its data, and RLHF data collection is a craft of its own.

\paragraph{The comparison format.} Annotators see a prompt and two or more responses sampled from
the current model, typically at temperature near 1 so they differ, and choose the better one. Some setups
collect a full ranking of $K$ responses (InstructGPT used $K$ between 4 and 9), a Likert-scale
strength (``much better'', ``slightly better''), or separate judgements per axis (helpfulness,
harmlessness, honesty). A ranking of $K$ responses yields $\binom{K}{2}$ pairwise comparisons for
the cost of reading $K$ responses, which is why rankings were popular.

\paragraph{Guidelines and agreement.} Labeling guidelines define what ``better'' means: following the
instruction, truthfulness, avoiding harm, and so on. Even with good guidelines, independent
annotators agree on only about three-quarters of comparisons in the published RLHF datasets.
That agreement rate is a rough ceiling on the held-out accuracy any reward model can reach.

\paragraph{Who labels matters.} The annotator pool's demographics, expertise and stylistic tastes
shift the resulting model. A strong, well-documented effect is \emph{length}: both human
and LLM judges prefer longer, more structured answers (bullet points, headings), and
policies learn to exploit this. Expert annotation for maths, code and law is now a large
industry. AI feedback (\cref{sec:rlhf:rlaif}) is the other answer to the cost and consistency problem.

\paragraph{Practical data plumbing.} Store preference data as JSONL with one comparison per line and
\emph{keep the provenance}: which model and which sampling settings produced each response, the annotator ID, and
the time. Split train and validation \emph{by prompt}, never by pair, so the validation accuracy
measures generalisation to new prompts. Collect comparisons on samples from the \emph{current}
policy in rounds. A reward model trained only on SFT outputs is out of distribution for a
policy that PPO has already moved, the shift behind the over-optimisation in
\cref{sec:rlhf:overopt}.

\begin{lstlisting}[style=plain,caption={One line of a preference dataset, with the metadata that later turns out to matter.},label={lst:rlhf:jsonl}]
{"prompt_id": "p-01733", "prompt": [{"role": "user", "content": "..."}],
 "chosen":   {"text": "...", "model": "sft-v3", "temperature": 1.0, "seed": 17},
 "rejected": {"text": "...", "model": "sft-v3", "temperature": 1.0, "seed": 18},
 "strength": "slightly_better", "axis": "helpfulness",
 "annotator": "a-112", "round": 2, "split": "train"}
\end{lstlisting}

\section{Reward models}\label{sec:rlhf:rm}

\subsection{The Bradley--Terry model as a latent-utility model}

Suppose every response $y$ to prompt $x$ has a latent quality $r^\star(x,y)$, and an annotator
perceives $r^\star + \epsilon$ with independent noise $\epsilon$. If the noise follows a Gumbel
distribution, the difference of two noises is logistic, and the probability that the annotator
prefers $y_w$ (``win'') over $y_l$ (``lose'') is
\begin{equation}
  \Pr[y_w \succ y_l \mid x] \;=\; \sigma\!\big(r^\star(x,y_w) - r^\star(x,y_l)\big),
  \qquad \sigma(z)=\frac{1}{1+e^{-z}}.
  \label{eq:rlhf:bt}
\end{equation}
This is the Bradley--Terry model from \cref{sec:eval:leaderboards}, with responses in place of
chatbots. Fitting a parametric $r_\phi$ by maximum likelihood gives the reward-model loss
\begin{equation}
  \mathcal{L}_{\text{RM}}(\phi) = -\E_{(x,y_w,y_l)}\,\log\sigma\!\big(r_\phi(x,y_w)-r_\phi(x,y_l)\big).
  \label{eq:rlhf:rmloss}
\end{equation}
Two properties follow directly. First, the loss depends only on \emph{differences}, so
$r_\phi$ is defined up to a per-prompt constant and the reward has no natural zero. Second, the
gradient with respect to the margin $\Delta = r_w - r_l$ is $-(1-\sigma(\Delta))$, so pairs the
model already ranks confidently stop contributing, which is the same self-weighting we will
see in DPO.

\paragraph{Margins and ties.} When annotators report \emph{how much} better, a margin
$m\ge 0$ can be subtracted inside the sigmoid, $-\log\sigma(r_w-r_l-m)$, so ``much better'' pairs
must be separated by more. Llama~2 (Touvron et al., 2023) trained its reward models this way.
Ties can be dropped, or modelled with a symmetric loss that pulls the two scores together.

\paragraph{$K$-way rankings.} With $K$ ranked responses per prompt, form all $\binom{K}{2}$ pairs.
InstructGPT's important engineering detail: put all pairs from one prompt \emph{in the same batch
element} and average the loss over the pairs of that prompt. If the pairs are shuffled into the
dataset independently, each response appears in $K-1$ separate gradient steps within an epoch and
the reward model overfits after one epoch. Grouping also saves compute: the $K$ responses each go
through one forward pass, not $K-1$.

\begin{lstlisting}[caption={Bradley--Terry and the grouped $K$-way ranking loss (\ident{llmsys/rlhf/losses.py}).},label={lst:rlhf:btloss}]
def bt_loss(r_chosen, r_rejected, margin=None):
    diff = r_chosen - r_rejected
    if margin is not None:
        diff = diff - margin
    return -F.logsigmoid(diff).mean()

def ranking_loss(scores, ranks):          # scores, ranks: (P, K); rank 0 = best
    P, K = scores.shape
    losses = []
    for i, j in itertools.combinations(range(K), 2):
        better_i = (ranks[:, i] < ranks[:, j]).float()
        diff = scores[:, i] - scores[:, j]
        oriented = better_i * diff + (1 - better_i) * (-diff)
        losses.append(-F.logsigmoid(oriented))
    # average within each prompt, then across prompts
    return torch.stack(losses, dim=1).mean()
\end{lstlisting}

\subsection{Architecture: a language model with a scalar head}

A reward model is the policy architecture with the vocabulary projection
$\R^{d}\to\R^{|V|}$ replaced by a scalar projection $\R^{d}\to\R$. The prompt and response are
concatenated and run through the trunk, and the hidden state of the \emph{last real token}
(normally the end-of-sequence token) is projected to a score. Because the model is causal, the
last position is the only one that has seen the whole response. Starting from the SFT model is
standard: the trunk already ``reads'' the domain, and the head is the only fresh layer.

\begin{figure}[ht]
\centering
\resizebox{\linewidth}{!}{%
\begin{tikzpicture}[tok/.style={draw, minimum width=7mm, minimum height=6mm, font=\scriptsize},
  >={Stealth[length=2mm]}]
\def\row#1#2#3#4{% y, label, n_resp, fill
  \node[font=\scriptsize, anchor=east] at (-0.5,#1) {#2};
  \foreach \i in {0,...,3} \node[tok, fill=fdata!18] at (0.75*\i,#1) {$x_{\i}$};
  \foreach \i in {1,...,#3} \node[tok, fill=#4] at (3+0.75*\i-0.75,#1) {$y_{\i}$};
}
\row{0}{chosen $y_w$}{5}{keygreen!18}
\node[tok, fill=keygreen!45] at (3+0.75*5,0) {\textsc{eos}};
\node[tok, fill=black!8] at (3+0.75*6,0) {pad};
\row{-0.9}{rejected $y_l$}{3}{warnred!15}
\node[tok, fill=warnred!40] at (3+0.75*3,-0.9) {\textsc{eos}};
\foreach \k in {4,5,6} \node[tok, fill=black!8] at (3+0.75*\k,-0.9) {pad};
\node[font=\scriptsize, text width=82mm, align=left, anchor=north west] at (-0.4,-1.4)
  {Rows $1..B$ hold the chosen responses and rows $B{+}1..2B$ the rejected ones, so a pair shares one
   batch. The score is read at each row's own last real token (gather by length).};
\node[draw, rounded corners, fill=fcompute!15, minimum width=34mm, minimum height=17mm,
  align=center, font=\small] (trunk) at (10.6,-0.45) {transformer trunk\\(SFT weights)\\[2pt]
  {\scriptsize one forward over $2B$ rows}};
\draw[->] (8.1,-0.45) -- (trunk.west);
\node[draw, rounded corners, fill=napkin!20, font=\small] (h1) at (14.2,0) {$\mathbf{w}^\top h_{\textsc{eos}}$};
\node[draw, rounded corners, fill=napkin!20, font=\small] (h2) at (14.2,-0.9) {$\mathbf{w}^\top h_{\textsc{eos}}$};
\draw[->] (trunk.east |- h1) -- (h1); \draw[->] (trunk.east |- h2) -- (h2);
\node[font=\small] (rw) at (15.9,0) {$r_w$}; \node[font=\small] (rl) at (15.9,-0.9) {$r_l$};
\draw[->] (h1)--(rw); \draw[->] (h2)--(rl);
\node[draw, rounded corners, fill=white, font=\small] (loss) at (18.4,-0.45) {$-\log\sigma(r_w-r_l)$};
\draw[->] (rw) -- (loss.north west); \draw[->] (rl) -- (loss.south west);
\end{tikzpicture}}
\caption{Reward-model training. The chosen and rejected responses of a pair are stacked into one
right-padded batch. The trunk runs once, the scalar head reads each row at its own
\textsc{eos} position, and the BT loss compares the two scores.}
\label{fig:rlhf:rmbatch}
\end{figure}

\begin{lstlisting}[caption={A reward model and a value model share one trunk class; they differ only in which positions they read (\ident{llmsys/rlhf/models.py}).},label={lst:rlhf:models}]
def trunk_hidden(model, ids, pad_mask=None):          # Qwen3 trunk, no LM head
    x = model.tok_emb(ids)
    for blk in model.trf_blocks:
        x = blk(x, model.cos, model.sin, pad_mask=pad_mask)
    return model.final_norm(x)

class RewardModel(ScalarHeadModel):
    def forward(self, ids, lengths):                     # -> (B,)
        h = self.hidden(ids)
        last = (lengths - 1).clamp(min=0)
        return self.head(h[torch.arange(ids.shape[0]), last]).squeeze(-1)

class ValueModel(ScalarHeadModel):
    def forward(self, ids):                       # -> (B, T): V(s_t) per prefix
        return self.head(self.hidden(ids)).squeeze(-1)
\end{lstlisting}

\begin{pitfall}[Where the score is read]
Right padding plus ``take position $-1$'' reads a \emph{pad} token for every row except the longest.
The model then learns a function of padding length, and its accuracy can look fine on a
validation set with the same length distribution. Always gather by length (or by the \textsc{eos}
index). With left padding, take position $-1$, but then the absolute positions differ from those
at generation time unless RoPE is used (\cref{sec:serving:masks}). Test it: shuffle the padding
amount and check that the scores do not change.
\end{pitfall}

\subsection{How big, how long, how good?}

\paragraph{Size.} InstructGPT used a 6B reward model for a 175B policy and reported that a 175B
reward model was less stable to train. Later pipelines more often use a reward model of the
policy's size or larger. The trade-off is systems-shaped: a larger reward model is a better judge, but it costs
a full forward pass over every rollout in every PPO iteration.

\paragraph{Epochs.} Reward models overfit quickly. One epoch over the comparisons is common, and
validation accuracy should be tracked every few hundred steps.

\paragraph{Normalisation.} Because BT fixes $r_\phi$ only up to a constant, many pipelines shift
the reward so that a reference set (for example SFT samples) scores zero on average, and some also
scale it to unit variance. This makes the KL coefficient $\beta$ transferable between reward models.

\paragraph{Evaluation.} The primary metric is \emph{pairwise accuracy} on held-out prompts. For
human-labelled data it is typically in the mid-60s to low-70s, close to inter-annotator agreement.
Check accuracy separately for each data source, response-length bucket and topic, because averages hide the
slices the policy will exploit. Benchmarks such as RewardBench (Lambert et al., 2024) compare reward models
on curated pairs. The real test is whether optimising against the reward model improves what you
care about, and only a downstream evaluation (\cref{ch:eval}) answers that.

\paragraph{Outcome versus process rewards.} The reward models here are \emph{outcome} reward models
(ORMs): one score for the whole response. A \emph{process} reward model (PRM) scores each step of a
reasoning chain (Lightman et al., 2023). That provides denser credit assignment, but step labels are costly
to collect. \Cref{ch:refine} used a PRM as a verifier at inference time. In PPO, a PRM's
per-step scores simply become additional per-token rewards in the shaped reward of
\cref{sec:rlhf:kl}.

\begin{systemslens}[Reward-model inference is pure prefill]
At PPO time the reward model never generates. It runs one prefill per response and reads one
position. That makes it compute-bound and easy to batch: sort rollouts by length to reduce
padding, run with large batches, and cache nothing. If the reward model and the policy share a
tokenizer and chat template, the policy's token IDs can be fed straight in. If not (a
different model family), every response must be detokenised and re-tokenised, a CPU cost that
belongs in the profile. Some systems serve the reward model as a separate service (for example a vLLM
instance in scoring mode) so it can scale independently of the trainer (\cref{sec:rlsys:reward}).
\end{systemslens}

\begin{lstlisting}[caption={Training the reward model on stacked pairs (\ident{llmsys/rlhf/ppo.py}, abridged).},label={lst:rlhf:rmtrain}]
def train_reward_model(rm, pairs, pad_id, epochs=3, batch_size=32, lr=1e-3, ...):
    opt = torch.optim.AdamW(rm.parameters(), lr=lr, weight_decay=0.0)
    for ep in range(epochs):
        rng.shuffle(pairs)
        for i in range(0, len(pairs), batch_size):
            batch = pairs[i: i + batch_size]
            ids, lens = pad_pairs(batch, pad_id)   # (2B, T)
            r = rm(ids, lens)
            n = len(batch)
            loss = loss_fn(r[:n], r[n:])                 # Bradley-Terry
            opt.zero_grad(set_to_none=True)
            loss.backward()
            torch.nn.utils.clip_grad_norm_(rm.parameters(), 1.0)
            opt.step()
\end{lstlisting}

\section{Per-token rewards and the KL penalty}\label{sec:rlhf:kl}

The reward model emits one number per response, but PPO assigns credit per token. RLHF bridges the two by
\emph{shaping} the reward: every generated token pays a KL toll, and the final token also collects the
reward-model score,
\begin{equation}
  r_t \;=\; -\beta\,\underbrace{\log\frac{\piold(y_t\mid s_t)}{\piref(y_t\mid s_t)}}_{\text{KL}_t\ (k_1\text{ estimator})}
  \;+\; \mathbb{1}[t = T]\; r_\phi(x, y),
  \qquad t = 1,\dots,T.
  \label{eq:rlhf:shaped}
\end{equation}
Summed over the response, the KL tolls are an unbiased single-sample estimate of
$\beta\,\KL(\piold\,\|\,\piref)$ for that prompt, so maximising $\sum_t r_t$ maximises the
objective~\eqref{eq:rlhf:objective}. The log-probabilities are those of the
policy \emph{at rollout time} ($\piold$); they are constants during the update.

\begin{figure}[ht]
\centering
\begin{tikzpicture}
\begin{groupplot}[group style={group size=2 by 1, horizontal sep=1.6cm},
  width=0.5\linewidth, height=4.6cm, xlabel={response token $t$},
  symbolic x coords={1,2,3,4,5,6,7}, xtick=data, axis lines*=left, enlarge x limits=0.08,
  tick label style={font=\scriptsize}, label style={font=\small}, title style={font=\small}]
\nextgroupplot[ybar, bar width=8pt, title={KL toll $-\beta\,\mathrm{KL}_t$ (every token)}, ymin=-0.14, ymax=0.04,
  ylabel={reward}, yticklabel style={/pgf/number format/fixed}]
\addplot[fill=warnred!45, draw=warnred] coordinates {(1,-0.02) (2,-0.06) (3,0.01) (4,-0.12) (5,-0.03) (6,-0.05) (7,-0.04)};
\nextgroupplot[ybar, bar width=8pt, title={total $r_t$ (score added at \textsc{eos})}, ymin=-0.2, ymax=2.2]
\addplot[fill=napkin!55, draw=napkin] coordinates {(1,-0.02) (2,-0.06) (3,0.01) (4,-0.12) (5,-0.03) (6,-0.05) (7,1.96)};
\end{groupplot}
\end{tikzpicture}
\caption{The shaped reward of one response (illustrative values, $\beta=0.1$, reward-model
score $2.0$). The KL toll (left, note the scale) is paid token by token and can be slightly positive where the policy is \emph{less} confident than
the reference; the reward-model score arrives only at the end-of-sequence token. GAE spreads that
terminal reward back over the response.}
\label{fig:rlhf:shaped}
\end{figure}

\paragraph{KL in the reward or in the loss?} PPO-RLHF puts the KL penalty \emph{inside the reward},
so it flows through the critic and GAE like any other reward and is subject to clipping. GRPO
(\cref{sec:grpo2:kl}) instead adds an explicit $\beta\,\widehat{\KL}$ term to the
\emph{loss}, usually with the non-negative $k_3$ estimator, because it has no critic to carry a per-token reward.
The two are not equivalent. In the reward, the penalty's gradient is a policy-gradient estimate of
$\nabla\KL$; in the loss, it is the gradient of the estimator, which is unbiased for $\nabla\KL$
only for particular estimator choices (\cref{ch:grpo2} discusses this subtlety). Our code
implements the $k_1$, $k_2$ and $k_3$ estimators so you can compare them in the lab.

\paragraph{Choosing $\beta$: fixed or adaptive.} With a fixed $\beta$, the KL the policy ends up
spending depends on how exploitable the reward model is. Ziegler et al.\ (2019) introduced an
\emph{adaptive} controller that targets a KL budget:
\begin{equation}
  \beta \leftarrow \beta\Big(1 + \operatorname{clip}\big(\tfrac{\mathrm{KL}-\mathrm{KL}_{\text{target}}}{\mathrm{KL}_{\text{target}}},
  -0.2,\,0.2\big)\cdot \tfrac{n}{H}\Big),
\end{equation}
where $n$ is the number of samples since the last update and $H$ is a horizon. This is a
proportional controller with a bounded gain: if the policy overspends, $\beta$ rises, and if it
underspends, $\beta$ falls. Targeting KL makes runs comparable across reward models with
different scales.

\paragraph{Missing \textsc{eos}.} A response truncated at the length limit has no well-defined
score, because the reward model was trained on complete responses. Common practice is to replace its
score with a fixed penalty, which also teaches the policy to finish. Our trainer does this with
\code{missing\_eos\_score}.

\begin{lstlisting}[caption={KL estimators and the shaped reward (\ident{llmsys/rlhf/losses.py}).},label={lst:rlhf:shaped}]
def kl_per_token(logp, logp_ref, kind="k1"):
    log_ratio = logp - logp_ref                # log(pi/pi_ref), sampled token
    if kind == "k1": return log_ratio                 # unbiased, can be negative
    if kind == "k2": return 0.5 * log_ratio ** 2      # biased, low variance
    if kind == "k3": return torch.expm1(-log_ratio) + log_ratio   # unbiased, >= 0

def kl_shaped_rewards(score, logp, logp_ref, mask, beta, kind="k1"):
    kl = kl_per_token(logp, logp_ref, kind) * mask
    rewards = -beta * kl
    last = (mask.sum(1).long() - 1).clamp(min=0)   # each row's final token
    rewards = rewards.clone()
    rewards[torch.arange(rewards.shape[0]), last] += score
    return rewards * mask, kl
\end{lstlisting}

\section{The value model}\label{sec:rlhf:value}

\paragraph{What it predicts.} $V_\psi(s_t)$ estimates the expected shaped return from state
$s_t = (x, y_{<t})$ onwards, $\E\big[\sum_{k\ge t} r_k\big]$ under the current policy. With
$\gamma=1$ and a terminal-only score, that is ``the score this response will end up with,
minus the KL still to be paid, judged from a prefix''. The value at the prompt's last token is
a per-prompt baseline: some prompts are easy to score well on and some are not. Values at later
tokens let the critic say \emph{which} tokens moved the expected outcome.

\paragraph{Where it reads.} Column $t$ of the value model's output values the state that
\emph{precedes} response token $t$. For a prompt of length $P$, the values aligned with response
tokens are therefore at positions $P-1,\dots,P+T-2$, exactly the same slice as the logits that
produce those tokens. Getting this off by one is the classic critic bug: the critic then sees the
token whose value it is predicting.

\paragraph{Initialisation and architecture.} InstructGPT initialised the value model from the
\emph{reward model}, which already maps text to a score. Other implementations start from the
SFT policy and add a fresh head. A separate network doubles the trained memory
(\cref{sec:rlhf:system}). Sharing the policy's trunk with a second head halves it, but couples
two losses with very different scales on one set of weights, which can destabilise both. Most
large-scale systems keep them separate. A common trick is a short \emph{critic warm-up}: train
only the value model for the first few iterations, so the first policy updates do not use
garbage advantages.

\paragraph{Its loss.} The value model regresses onto the GAE returns of the next section, with
PPO2-style clipping that forbids moving more than $\epsilon_v$ away from the rollout-time
prediction in one update:
\begin{equation}
  \mathcal{L}_V(\psi) = \tfrac12\,\E_t\Big[\max\big((V_\psi(s_t)-\hat R_t)^2,\;
   (V^{\text{old}}_t + \clip(V_\psi(s_t)-V^{\text{old}}_t,-\epsilon_v,\epsilon_v)-\hat R_t)^2\big)\Big].
  \label{eq:rlhf:vloss}
\end{equation}
Value clipping is controversial: it was introduced as an implementation detail and some studies
find it neutral or harmful. Keep it configurable (\code{vf\_clip=None} disables it).

\paragraph{Degenerate cases worth knowing.} If $\lambda=1$ and $\gamma=1$, GAE ignores the critic
except as a baseline at each token, so a bad critic only adds variance, not bias. If $\lambda=0$, the
advantage is the one-step TD error $\delta_t$, which relies entirely on the critic, so a bad critic
gives biased advantages. If the critic predicts a constant, PPO degenerates to REINFORCE with a
constant baseline. Measure the critic with \textbf{explained variance},
$1-\mathrm{Var}[\hat R-V]/\mathrm{Var}[\hat R]$: near 0 means the critic is useless, and negative means it
is actively harmful.

\section{Generalised advantage estimation}\label{sec:rlhf:gae}

Given per-token rewards and values, how much credit should token $t$ get? Define the TD residual
\begin{equation}
  \delta_t = r_t + \gamma V(s_{t+1}) - V(s_t), \qquad V(s_{T+1}) := 0 \text{ (terminal)}.
\end{equation}
$\delta_t$ is a one-step advantage estimate: what happened at step $t$, plus what the critic now
expects, minus what it expected before. Summing $k$ residuals gives the $k$-step estimate
$\hat A^{(k)}_t = \sum_{l=0}^{k-1}\gamma^l\delta_{t+l}$, which telescopes to
$\sum_{l<k}\gamma^l r_{t+l} + \gamma^k V(s_{t+k}) - V(s_t)$. Small $k$ trusts the critic (low
variance, biased if the critic is wrong); large $k$ trusts the sampled rewards (unbiased, high
variance). GAE (Schulman et al., 2016) averages all $k$ with geometric weights $(1-\lambda)\lambda^{k-1}$,
which collapses to a simple recursion:
\begin{equation}
  \hat A_t \;=\; \sum_{l\ge0} (\gamma\lambda)^l\,\delta_{t+l}
  \;=\; \delta_t + \gamma\lambda\,\hat A_{t+1},
  \qquad \hat R_t = \hat A_t + V(s_t).
  \label{eq:rlhf:gae}
\end{equation}
The returns $\hat R_t$ are the value model's regression targets. $\lambda=0$ gives the TD error, and
$\lambda=1$ gives the Monte Carlo return minus the baseline. For LLMs $\gamma=1$ and
$\lambda\in[0.9,1]$ are typical.

\begin{table}[ht]
\centering\small
\begin{tabular}{@{}lcccc@{}}
\toprule
token $t$ & 1 & 2 & 3 & 4 (\textsc{eos})\\
\midrule
$\mathrm{KL}_t=\log\piold-\log\piref$ & $0.20$ & $0.00$ & $-0.50$ & $0.10$\\
$r_t$ ($\beta=0.1$, score $2.0$) & $-0.02$ & $0.00$ & $0.05$ & $1.99$\\
$V(s_t)$ & $0.50$ & $0.80$ & $1.20$ & $1.60$\\
$\delta_t$ & $0.28$ & $0.40$ & $0.45$ & $0.39$\\
\midrule
$\hat A_t$, $\lambda=0$ & $0.28$ & $0.40$ & $0.45$ & $0.39$\\
$\hat A_t$, $\lambda=0.95$ & $1.40$ & $1.18$ & $0.82$ & $0.39$\\
$\hat A_t$, $\lambda=1$ & $1.52$ & $1.24$ & $0.84$ & $0.39$\\
$\hat R_t$, $\lambda=0.95$ & $1.90$ & $1.98$ & $2.02$ & $1.99$\\
\bottomrule
\end{tabular}
\caption{GAE on a four-token response ($\gamma=1$), computed with \code{gae()}. With $\lambda=1$ the
first token's advantage is the total reward minus the first value, $2.02-0.50$. With $\lambda=0$
every token sees only its own TD error.}
\label{tab:rlhf:gae}
\end{table}

\begin{figure}[ht]
\centering
\resizebox{0.95\linewidth}{!}{%
\begin{tikzpicture}[tok/.style={draw, rounded corners=2pt, minimum width=13mm, minimum height=8mm,
  font=\small, fill=fcompute!12}, >={Stealth[length=2mm]}]
\foreach \i/\lbl in {0/$y_1$,1/$y_2$,2/$y_3$,3/$y_4$,4/\textsc{eos}}
  \node[tok] (t\i) at (2.2*\i,0) {\lbl};
\foreach \i in {0,...,4} \node[font=\scriptsize] at (2.2*\i,-0.75) {$\delta_{\the\numexpr\i+1\relax}$};
\node[draw, rounded corners, fill=napkin!25, font=\small] (rm) at (11.6,0) {$r_\phi$};
\draw[->] (rm) -- (t4);
\foreach \i in {4,3,2,1} {
  \pgfmathtruncatemacro{\j}{\i-1}
  \draw[->, thick, fcompute] (t\i.north) to[out=140, in=40]
    node[above, font=\scriptsize] {$\times\gamma\lambda$} (t\j.north);
}
\end{tikzpicture}}
\caption{GAE as a backward scan, $\hat A_t=\delta_t+\gamma\lambda\hat A_{t+1}$: one right-to-left
pass per response, with padding masked so the recursion stops at each row's own \textsc{eos}.
Each token's advantage is its own TD residual plus a discounted copy of its successor's advantage, so the terminal reward-model score flows backwards, attenuated by
$\gamma\lambda$ per step and corrected by the critic's predictions along the way.}
\label{fig:rlhf:gae}
\end{figure}

\begin{lstlisting}[caption={GAE over right-padded responses (\ident{llmsys/rlhf/losses.py}).},label={lst:rlhf:gae}]
def gae(rewards, values, mask, gamma=1.0, lam=0.95):
    B, T = rewards.shape
    adv = torch.zeros_like(rewards)
    last_gae = torch.zeros(B)
    for t in reversed(range(T)):
        nv = values[:, t + 1] * mask[:, t + 1] if t + 1 < T else torch.zeros(B)
        delta = rewards[:, t] + gamma * nv - values[:, t]
        carry = mask[:, t + 1] if t + 1 < T else 0.0   # stop after final token
        last_gae = (delta + gamma * lam * last_gae * carry) * mask[:, t]
        adv[:, t] = last_gae
    return adv * mask, (adv + values) * mask             # advantages, returns
\end{lstlisting}

\begin{systemslens}[GAE is cheap, but sequential]
GAE is $O(BT)$ arithmetic, negligible next to the forward passes, but the scan over $t$ is
sequential. In PyTorch a Python loop over $T$ launches a few kernels per step: fine for $T$ in the
hundreds, noticeable for 16k-token reasoning traces on a GPU. Options: run it on the CPU in NumPy (the
tensors are small), write the scan as a reversed cumulative sum when $\gamma\lambda$ is constant and
there is no masking inside a row, or fuse it into one Triton kernel with one program per row
(\cref{app:triton}). Whiten advantages \emph{across the whole batch} (masked mean and variance over
all real tokens), not per sequence: per-sequence whitening removes the signal that one response was
better than another.
\end{systemslens}

\section{The PPO objective for RLHF}\label{sec:rlhf:ppo}

With advantages in hand, the policy update is the clipped surrogate from \cref{ch:pg} and
\cref{sec:grpo2:clip}, applied per token:
\begin{equation}
  \mathcal{L}_\pi(\theta) = -\,\E_t\Big[\min\big(\rho_t\hat A_t,\;
     \clip(\rho_t, 1-\epsilon, 1+\epsilon)\hat A_t\big)\Big],
  \qquad \rho_t = \frac{\pitheta(y_t\mid s_t)}{\piold(y_t\mid s_t)}.
  \label{eq:rlhf:clip}
\end{equation}
The expectation is a masked mean over all real response tokens in the minibatch. The full PPO
objective adds the critic and an optional entropy bonus:
\begin{equation}
  \mathcal{L}(\theta,\psi) = \mathcal{L}_\pi(\theta) + c_v\,\mathcal{L}_V(\psi) - c_H\,\E_t\big[H(\pitheta(\cdot\mid s_t))\big]
  \;\big[\, + \gamma_{\text{ptx}}\,\mathcal{L}_{\text{LM}}(\theta)\,\big].
  \label{eq:rlhf:total}
\end{equation}
When actor and critic are separate networks with separate optimisers, as in our trainer, $c_v$
only matters if they share parameters. The entropy bonus is usually zero for LLMs, because the KL
penalty already keeps the policy from collapsing.

\paragraph{PPO-Clip versus PPO-Penalty.} The original PPO paper offered two ways to keep the new policy near
the old one: clip the ratio, or penalise $\KL(\piold\|\pitheta)$ with an adaptive coefficient. The
clipped version won in practice and is what ``PPO'' means today. Do not confuse it with the
KL-to-\emph{reference} penalty of RLHF. PPO-RLHF uses \emph{both} a trust region (clipping,
relative to $\piold$, which changes every iteration) and an anchor (KL in the reward, relative to
$\piref$, which never changes).

\paragraph{PPO-ptx.} InstructGPT observed an \emph{alignment tax}: RLHF lowered performance on some
public NLP benchmarks. Their fix mixes a pretraining language-modelling loss into every PPO update
(the bracketed term in \cref{eq:rlhf:total}), so the policy keeps its general capabilities. The cost is an extra
forward and backward pass on a pretraining batch per minibatch. Our trainer exposes it as
\code{ptx\_coef} with a pluggable batch function.

\begin{lstlisting}[caption={Clipped policy loss and clipped value loss (\ident{llmsys/rlhf/losses.py}).},label={lst:rlhf:losses}]
def ppo_policy_loss(logp, logp_old, adv, mask, clip=0.2):
    log_ratio = (logp - logp_old) * mask
    ratio = torch.exp(log_ratio)
    clipped = torch.clamp(ratio, 1 - clip, 1 + clip)
    surr = torch.minimum(ratio * adv, clipped * adv)
    loss = masked_mean(-surr, mask)
    outside = ((ratio - 1).abs() > clip).float()
    stats = {"clipfrac": float(masked_mean(outside, mask)),
             "approx_kl": float(masked_mean(0.5 * log_ratio ** 2, mask))}
    return loss, stats

def value_loss(values, values_old, returns, mask, clip=0.2):
    err = (values - returns) ** 2
    if clip is not None:
        v_clip = values_old + torch.clamp(values - values_old, -clip, clip)
        err = torch.maximum(err, (v_clip - returns) ** 2)
    return 0.5 * masked_mean(err, mask)
\end{lstlisting}

\begin{pitfall}[Token mean, sequence mean, and length]
\code{masked\_mean} over all tokens of a minibatch weights long responses more than short ones,
since they contribute more terms. A per-sequence mean followed by a batch mean does the opposite.
Neither is ``correct''. The choice changes the effective objective, and with it the length
dynamics (\cref{sec:grpo2:more} catalogues the variants and \cref{ch:pg} derives the bias).
Pick one deliberately, log the mean response length, and keep the choice fixed when comparing
runs.
\end{pitfall}

\section{The PPO-RLHF loop}\label{sec:rlhf:loop}

We now have every component. One PPO iteration has three phases and two nested loops: an outer
loop that collects fresh rollouts, and an inner loop that makes several optimisation passes over them.

\begin{table}[ht]
\centering\small
\begin{tabularx}{\linewidth}{@{}lccXX@{}}
\toprule
\textbf{Model} & \textbf{Trained?} & \textbf{Initialised from} & \textbf{Input $\to$ output} & \textbf{Role}\\
\midrule
Policy $\pitheta$ & yes & SFT model & prefix $\to$ next-token distribution & generates rollouts; the model we ship\\
Reference $\piref$ & no & SFT model & prefix $\to$ log-prob of the sampled token & anchors the policy through the KL toll\\
Reward $r_\phi$ & no (here) & SFT trunk + scalar head & (prompt, response) $\to$ one scalar & stands in for the human judge\\
Value $V_\psi$ & yes & reward model or SFT trunk & every prefix $\to$ one scalar & per-token baseline for GAE\\
\bottomrule
\end{tabularx}
\caption{The four models of PPO-RLHF. Only two are trained, but all four run a forward pass on every
rollout, and the policy additionally generates them.}
\label{tab:rlhf:roles}
\end{table}

\begin{lstlisting}[style=plain,caption={One PPO-RLHF iteration in pseudocode.},label={lst:rlhf:pseudo}]
repeat for iteration = 1, 2, ...:
  # Phase 1: rollout (policy in inference mode, KV cache)
  sample B prompts x;  y ~ pi_theta(. | x)       -> responses, mask, has_eos
  # Phase 2: score (four forward passes, no gradients)
  logp_old = log pi_theta(y_t | s_t) logp_ref = log pi_ref(y_t | s_t)
  V_old    = V_psi(s_t)              score    = r_phi(x, y)  (penalty if no eos)
  r_t      = -beta * (logp_old - logp_ref) + [t = T] * score
  A_t, R_t = GAE(r, V_old; gamma, lambda);  A_t = whiten(A_t over the batch)
  # Phase 3: learn (ppo_epochs passes, shuffled minibatches)
  for epoch in 1..E:  for minibatch in shuffle(B) / M:
      L_pi = clipped surrogate(logp_new / logp_old, A)   -> step actor
      L_V  = clipped value loss(V_new, V_old, R)         -> step critic
  update beta with the KL controller
\end{lstlisting}

Three subtleties trip up first implementations:

\begin{itemize}
  \item \textbf{``Old'' means rollout-time.} \code{logp\_old} and \code{V\_old} are computed once in
        phase 2 and frozen. During phase 3 the ratio starts at exactly 1 on the first minibatch,
        and drifts as the epochs proceed; clipping only becomes active after the first update.
        Recomputing \code{logp\_old} every epoch silently turns PPO into a sequence of tiny
        on-policy steps without a trust region.
  \item \textbf{Sampling and scoring must agree.} The log-probabilities used in the ratio must be those
        of the distribution you \emph{sampled from}. If generation used temperature $\tau$ or top-$p$,
        either score with the same transformation or accept an off-policy bias. In large systems the
        rollout engine and the trainer compute slightly different numbers even with identical settings,
        which is the train/inference mismatch of \cref{sec:rlsys:mismatch}.
  \item \textbf{Frozen means frozen.} The reference and reward models must never receive gradients, never
        be put in train mode (dropout), and never be updated by an optimizer that was built over all
        parameters. Freeze them explicitly and assert it.
\end{itemize}

\begin{figure}[ht]
\centering
\resizebox{\linewidth}{!}{%
\begin{tikzpicture}[
  ph/.style={draw, rounded corners=2pt, minimum height=8mm, font=\scriptsize, align=center},
  >={Stealth[length=2mm]}]
\node[font=\scriptsize\bfseries, anchor=east] at (-0.2,1.6) {policy $\pi_\theta$};
\node[font=\scriptsize\bfseries, anchor=east] at (-0.2,0.8) {reference $\pi_{\text{ref}}$};
\node[font=\scriptsize\bfseries, anchor=east] at (-0.2,0.0) {reward $r_\phi$};
\node[font=\scriptsize\bfseries, anchor=east] at (-0.2,-0.8) {value $V_\psi$};
% rollout
\node[ph, fill=fmemory!30, minimum width=50mm, anchor=west] at (0,1.6) {generate (decode, memory-bound)};
% score
\node[ph, fill=fcompute!20, minimum width=11mm, anchor=west] at (5.1,1.6) {fwd};
\node[ph, fill=fcompute!20, minimum width=11mm, anchor=west] at (5.1,0.8) {fwd};
\node[ph, fill=fcompute!20, minimum width=11mm, anchor=west] at (5.1,0.0) {fwd};
\node[ph, fill=fcompute!20, minimum width=11mm, anchor=west] at (5.1,-0.8) {fwd};
\node[ph, fill=white, minimum width=10mm, anchor=west] at (6.3,0.4) {GAE};
% learn
\foreach \e in {0,1,2,3} {
  \node[ph, fill=fcompute!40, minimum width=13mm, anchor=west] at (7.5+1.4*\e,1.6) {fwd+bwd};
  \node[ph, fill=fcompute!40, minimum width=13mm, anchor=west] at (7.5+1.4*\e,-0.8) {fwd+bwd};
}
\draw[decorate, decoration={brace, amplitude=4pt}] (0,2.15) -- node[above=4pt, font=\scriptsize]{phase 1: rollout} (5.0,2.15);
\draw[decorate, decoration={brace, amplitude=4pt}] (5.1,2.15) -- node[above=4pt, font=\scriptsize]{phase 2: score} (7.3,2.15);
\draw[decorate, decoration={brace, amplitude=4pt}] (7.5,2.15) -- node[above=4pt, font=\scriptsize]{phase 3: learn ($E$ epochs $\times$ minibatches)} (13.1,2.15);
\draw[->] (0,-1.5) -- node[below, font=\scriptsize]{time within one iteration} (13.1,-1.5);
\end{tikzpicture}}
\caption{Anatomy of one PPO-RLHF iteration. Only the policy and the value model are trained; the
reference and reward models run a single forward pass each. In practice the generation bar is often
the longest one, despite doing the fewest FLOPs, because decoding is memory-bound
(\cref{ch:inference}).}
\label{fig:rlhf:anatomy}
\end{figure}

\begin{lstlisting}[caption={The trainer's phases (\ident{llmsys/rlhf/ppo.py}, abridged). The losses come from an injectable module \code{fns}, so Lab 35 runs the whole system on your own implementations.},label={lst:rlhf:trainer}]
@torch.no_grad()
def score(self, ro):
    P, full = ro.prompts.shape[1], ro.full
    ro.logp_old   = response_logprobs(self.policy, full, P) * ro.mask
    ro.logp_ref   = response_logprobs(self.ref, full, P) * ro.mask
    # value of the state *before* each response token y_t
    ro.values_old = self.value(full)[:, P - 1: -1] * ro.mask
    score = self.rm(full, P + ro.mask.sum(1).long())
    penalty = torch.full_like(score, cfg.missing_eos_score)
    score = torch.where(ro.has_eos, score, penalty)
    ro.rewards, ro.kl = self.fns.kl_shaped_rewards(
        score, ro.logp_old, ro.logp_ref, ro.mask, self.kl_ctl.beta, cfg.kl_kind)
    adv, ro.returns = self.fns.gae(ro.rewards, ro.values_old, ro.mask,
                                   cfg.gamma, cfg.lam)
    ro.advantages = self.fns.whiten(adv, ro.mask)
    return ro

def learn(self, ro):
    for _ in range(cfg.ppo_epochs):
        for mb in torch.randperm(B).split(cfg.minibatch_size):
            m = ro.mask[mb]
            logp, ent = response_logprobs(self.policy, full[mb], P,
                                          with_entropy=True)
            pl, st = self.fns.ppo_policy_loss(logp * m, ro.logp_old[mb],
                                              ro.advantages[mb], m, cfg.clip)
            step(self.opt_pi, pl - cfg.ent_coef * masked_mean(ent, m))
            v = self.value(full[mb])[:, P - 1: -1] * m
            vl, _ = self.fns.value_loss(v, ro.values_old[mb], ro.returns[mb],
                                        m, cfg.vf_clip)
            step(self.opt_v, vl)
\end{lstlisting}

\section{Setting it up end to end}\label{sec:rlhf:setup}

RLHF is a pipeline of four jobs, and each one produces an artifact the next one consumes. Treat it
like a build system: every stage writes a versioned checkpoint plus a small JSON manifest saying
what it was built from.

\begin{lstlisting}[style=plain,caption={A project layout for an RLHF run.},label={lst:rlhf:layout}]
rlhf-run/
  configs/  sft.yaml  rm.yaml  ppo.yaml     # one config per stage, in git
  data/     sft.jsonl  prefs_r1.jsonl  prefs_r2.jsonl  prompts_rl.jsonl
  ckpt/     sft/  rm-r1/  ppo-r1/  rm-r2/  ppo-r2/   # + manifest.json
  logs/     ppo-r1/metrics.jsonl  samples/iter_000050.jsonl
  eval/     rm_acc.json  winrate_vs_sft.json  safety.json
\end{lstlisting}

\begin{lstlisting}[style=plain,caption={A PPO configuration with the knobs that matter, and values that are reasonable defaults for a first 1B--8B run.},label={lst:rlhf:yaml}]
policy_init:  ckpt/sft            # trained copy
reference:    ckpt/sft            # frozen copy of the same weights
reward_model: ckpt/rm-r1
value_init:   ckpt/rm-r1          # InstructGPT-style critic init
rollout: {prompts_per_iter: 512, max_new_tokens: 512,
          temperature: 1.0, top_p: 1.0}
ppo:     {epochs: 2, minibatch_size: 64, clip: 0.2, vf_clip: 0.2,
          gamma: 1.0, lam: 0.95, lr_policy: 1.0e-6, lr_value: 5.0e-6,
          max_grad_norm: 1.0, ent_coef: 0.0}
kl:      {init_beta: 0.05, target: 6.0, horizon: 10000,
          estimator: k1}          # target in nats per response
reward:  {normalize: true, missing_eos_score: -1.0, clip: 10.0}
ptx:     {coef: 0.0, data: data/pretrain_sample.bin}
log:     {every: 1, samples_every: 25, eval_every: 100}
\end{lstlisting}

The commands follow the stages. In the lab they are functions, but the shape is the same
at scale:

\begin{lstlisting}[style=shell,caption={The four stages as commands (the lab's \code{run.py} mirrors this).},label={lst:rlhf:cmds}]
python run.py sft   --config configs/sft.yaml --out ckpt/sft
python run.py prefs --policy ckpt/sft --n 4096 \
                    --out data/prefs_r1.jsonl        # sample + label pairs
python run.py rm    --config configs/rm.yaml  --out ckpt/rm-r1   # BT training
python run.py ppo   --config configs/ppo.yaml --out ckpt/ppo-r1  # the loop
python run.py eval  --policy ckpt/ppo-r1 --baseline ckpt/sft     # win rate, KL
\end{lstlisting}

\begin{table}[ht]
\centering\small
\begin{tabularx}{\linewidth}{@{}lXX@{}}
\toprule
\textbf{Metric} & \textbf{Healthy behaviour} & \textbf{If not, suspect}\\
\midrule
ratio at first minibatch & exactly 1.0, clipfrac 0 & sampling and scoring disagree (temperature, mask offset, dropout on)\\
KL to reference & starts at 0, grows smoothly to the target & starts $>0$: ref is not a copy of the init policy\\
RM score & rises, then plateaus & flat: lr too low or advantages all near 0; erratic: lr too high\\
clipfrac & $0.05$--$0.3$ & $\approx 0$: updates too small; $>0.4$: lr or epochs too large\\
approx\_kl ($\piold\to\pitheta$) & small, $\sim10^{-3}$--$10^{-2}$ per iteration & large: steps too big for the trust region\\
value explained variance & rises from $\le0$ towards $0.5$--$0.9$ & stays $<0$: critic misaligned (off-by-one) or lr too low\\
response length & changes slowly & runaway growth: length-biased RM or loss aggregation\\
entropy & declines slowly & collapses fast: KL too weak, lr too high\\
\bottomrule
\end{tabularx}
\caption{The dashboard of a PPO-RLHF run and what deviations usually mean.}
\label{tab:rlhf:dashboard}
\end{table}

\begin{deepdive}[Sanity checks before spending GPU hours]
\begin{enumerate}
  \item \textbf{Zero-update test:} run phase 3 with learning rate 0. The ratio must be exactly 1 and the KL
        exactly 0 every iteration.
  \item \textbf{Reward-model probes:} hand-write five pairs where the answer is obvious and check the
        ordering. Then check that padding and batch composition do not change any score.
  \item \textbf{Oracle reward:} replace $r_\phi$ with a simple programmatic reward (for example ``contains
        the word \emph{thanks}''). PPO should learn it within a few dozen iterations; if not, the bug
        is in the RL code, not the reward model.
  \item \textbf{Critic-only warm-up:} freeze the policy, train the value model for a few iterations, and
        watch explained variance rise.
  \item \textbf{Read samples:} log a few prompts and responses every $N$ iterations, and read them. Reward
        hacking is usually visible in the text long before it shows up in any metric.
\end{enumerate}
\end{deepdive}

\section{Running it: PPO on a synthetic preference task}\label{sec:rlhf:results}

The lab's \code{llmsys.rlhf.task} defines a small world where ``human'' preferences are known
exactly, so we can measure what the reward model and PPO do to \emph{true} quality. A prompt
names one of four topics. A response is a short sequence of words (positive, neutral and negative words)
ending in \textsc{eos}. The hidden preference function rewards mentioning the topic's own
word, using distinct positive words, avoiding negative words, ending near six words, finishing with
\textsc{eos}, and \emph{not repeating words}. Simulated annotators compare two responses with
Bradley--Terry noise on that hidden score. The SFT ``base model'' is fluent but has no taste: it emits
2--9 distinct random words and stops.

The pipeline is exactly the one of \cref{lst:rlhf:cmds}, at toy scale: a 2-layer, 48-wide
Qwen3-shaped policy with untied embeddings; SFT for 150 steps; 1,024 comparisons sampled from
the SFT policy; a reward model initialised from the SFT trunk and trained for four epochs; then
PPO with a value model initialised from the SFT trunk.

The reward model reaches about $70\%$ held-out pairwise accuracy after four epochs, roughly where
real reward models sit relative to annotator agreement. Its labels came from SFT samples, which almost
never repeat a word, so it has \emph{no information} about the true preference function's repetition
penalty. \Cref{fig:rlhf:curves} shows PPO for three KL coefficients, starting from the same SFT model and
reward model.

\begin{figure}[t]
\centering
\begin{tikzpicture}
\begin{groupplot}[group style={group size=3 by 1, horizontal sep=1.25cm},
  width=0.37\linewidth, height=4.6cm, xlabel={PPO iteration}, xmin=0, xmax=71,
  tick label style={font=\scriptsize}, label style={font=\small}, title style={font=\small},
  every axis plot/.append style={line width=0.8pt, mark=none}, grid=major, grid style={black!8}]
\nextgroupplot[title={reward-model score (proxy)}]
\addplot[color=sysblue, draw=sysblue] coordinates {(1,1.188) (3,2.416) (5,3.095) (7,3.373) (9,3.927) (11,3.912) (13,3.904) (15,4.002) (17,3.945) (19,4.027) (21,3.950) (23,3.973) (25,4.086) (27,4.171) (29,4.094) (31,3.943) (33,4.042) (35,4.097) (37,4.117) (39,4.026) (41,4.107) (43,4.017) (45,3.995) (47,3.996) (49,4.079) (51,4.002) (53,4.137) (55,3.883) (57,4.097) (59,4.083) (61,3.949) (63,4.230) (65,4.082) (67,4.077) (69,4.196) (70,3.945)};
\addplot[color=keygreen, draw=keygreen] coordinates {(1,1.188) (3,2.513) (5,3.233) (7,3.645) (9,4.021) (11,4.242) (13,4.344) (15,4.373) (17,4.355) (19,4.290) (21,4.546) (23,4.401) (25,4.566) (27,4.584) (29,4.568) (31,4.535) (33,4.623) (35,4.622) (37,4.634) (39,4.624) (41,4.598) (43,4.604) (45,4.557) (47,4.606) (49,4.599) (51,4.628) (53,4.608) (55,4.633) (57,4.565) (59,4.659) (61,4.639) (63,4.650) (65,4.670) (67,4.656) (69,4.639) (70,4.657)};
\addplot[color=warnred, draw=warnred] coordinates {(1,1.188) (3,2.541) (5,3.264) (7,3.639) (9,3.818) (11,4.212) (13,4.407) (15,4.510) (17,4.614) (19,4.604) (21,4.704) (23,4.705) (25,4.735) (27,4.773) (29,4.781) (31,4.796) (33,4.779) (35,4.809) (37,4.761) (39,4.674) (41,4.777) (43,4.802) (45,4.801) (47,4.766) (49,4.788) (51,4.795) (53,4.812) (55,4.529) (57,4.800) (59,4.810) (61,4.825) (63,4.821) (65,4.816) (67,4.829) (69,4.821) (70,4.818)};

\nextgroupplot[title={true (hidden) score}]
\addplot[color=sysblue, draw=sysblue] coordinates {(1,-0.048) (3,0.973) (5,1.973) (7,2.355) (9,3.122) (11,3.278) (13,2.880) (15,3.389) (17,3.061) (19,3.133) (21,2.858) (23,3.209) (25,3.350) (27,3.798) (29,3.291) (31,3.192) (33,2.880) (35,3.327) (37,3.542) (39,3.192) (41,3.327) (43,3.178) (45,3.145) (47,3.383) (49,3.620) (51,3.123) (53,3.164) (55,2.950) (57,3.228) (59,3.211) (61,3.017) (63,3.567) (65,3.281) (67,3.233) (69,3.730) (70,3.434)};
\addplot[color=keygreen, draw=keygreen] coordinates {(1,-0.048) (3,1.055) (5,2.153) (7,2.488) (9,3.261) (11,3.827) (13,3.488) (15,3.952) (17,3.403) (19,3.923) (21,3.866) (23,3.778) (25,3.908) (27,4.030) (29,3.808) (31,4.159) (33,4.358) (35,4.189) (37,4.200) (39,4.531) (41,4.275) (43,4.414) (45,4.408) (47,4.502) (49,4.675) (51,4.442) (53,4.588) (55,4.591) (57,4.586) (59,4.242) (61,4.358) (63,4.405) (65,4.319) (67,4.638) (69,4.487) (70,4.595)};
\addplot[color=warnred, draw=warnred] coordinates {(1,-0.048) (3,1.044) (5,2.095) (7,2.845) (9,2.867) (11,3.614) (13,3.498) (15,3.280) (17,2.891) (19,2.392) (21,2.345) (23,1.803) (25,2.077) (27,2.055) (29,1.767) (31,1.859) (33,2.152) (35,2.247) (37,2.159) (39,1.678) (41,2.367) (43,2.513) (45,2.598) (47,2.780) (49,2.720) (51,2.847) (53,2.806) (55,2.517) (57,2.809) (59,2.763) (61,2.825) (63,2.895) (65,2.842) (67,2.872) (69,2.817) (70,2.778)};
\nextgroupplot[title={KL to reference (nats)}]
\addplot[color=sysblue, draw=sysblue] coordinates {(1,0.000) (3,0.667) (5,1.510) (7,2.180) (9,3.668) (11,3.293) (13,3.440) (15,3.741) (17,3.386) (19,3.772) (21,3.507) (23,3.459) (25,3.700) (27,4.079) (29,3.859) (31,3.344) (33,3.445) (35,3.971) (37,3.583) (39,3.679) (41,3.685) (43,3.300) (45,3.263) (47,3.778) (49,3.503) (51,3.810) (53,3.663) (55,3.266) (57,3.579) (59,3.860) (61,3.546) (63,4.364) (65,3.537) (67,3.538) (69,4.072) (70,3.459)};
\addplot[color=keygreen, draw=keygreen] coordinates {(1,0.000) (3,0.734) (5,1.839) (7,3.971) (9,4.959) (11,6.223) (13,7.008) (15,6.598) (17,6.575) (19,6.932) (21,7.969) (23,8.065) (25,8.717) (27,8.557) (29,8.902) (31,8.177) (33,9.334) (35,9.261) (37,9.199) (39,8.782) (41,8.768) (43,8.748) (45,7.567) (47,8.610) (49,8.303) (51,9.294) (53,9.647) (55,8.981) (57,8.809) (59,9.840) (61,9.030) (63,9.340) (65,9.537) (67,8.991) (69,9.065) (70,9.232)};
\addplot[color=warnred, draw=warnred] coordinates {(1,0.000) (3,0.732) (5,1.953) (7,3.969) (9,4.928) (11,7.452) (13,9.868) (15,12.151) (17,14.199) (19,15.876) (21,18.551) (23,19.266) (25,19.523) (27,20.122) (29,21.653) (31,22.022) (33,20.865) (35,20.877) (37,22.752) (39,24.375) (41,22.244) (43,22.322) (45,22.035) (47,21.788) (49,22.219) (51,21.994) (53,22.754) (55,21.831) (57,22.708) (59,23.319) (61,23.361) (63,23.222) (65,23.237) (67,23.046) (69,23.352) (70,23.496)};
\end{groupplot}
\end{tikzpicture}
\caption{PPO-RLHF on the synthetic preference task, measured on CPU in this environment (one seed;
the three runs share the same SFT model and reward model). All three runs raise the learned reward. With
$\beta=0.05$ the true score rises with it. With $\beta=0$ the policy spends over 20 nats of KL: the proxy
keeps climbing while the true score peaks at iteration 11 and then falls. With $\beta=0.2$ the policy
stays close to the reference and improves less.}
\label{fig:rlhf:curves}
\end{figure}

\begin{figure}[ht]
\centering
\begin{tikzpicture}
\begin{axis}[width=0.66\linewidth, height=5.4cm, xlabel={$\sqrt{\mathrm{KL}}$ to reference (nats$^{1/2}$)},
  ylabel={score}, grid=major, grid style={black!8}, tick label style={font=\scriptsize},
  legend style={font=\scriptsize, at={(1.03,0.5)}, anchor=west}, legend cell align=left]
\addplot[color=sysblue, only marks, mark=o, mark size=1.1pt] coordinates {(0.000,1.188) (0.817,2.416) (1.229,3.095) (1.477,3.373) (1.915,3.927) (1.815,3.912) (1.855,3.904) (1.934,4.002) (1.840,3.945) (1.942,4.027) (1.873,3.950) (1.860,3.973) (1.923,4.086) (2.020,4.171) (1.964,4.094) (1.829,3.943) (1.856,4.042) (1.993,4.097) (1.893,4.117) (1.918,4.026) (1.920,4.107) (1.817,4.017) (1.806,3.995) (1.944,3.996) (1.872,4.079) (1.952,4.002) (1.914,4.137) (1.807,3.883) (1.892,4.097) (1.965,4.083) (1.883,3.949) (2.089,4.230) (1.881,4.082) (1.881,4.077) (2.018,4.196)};
\addplot[color=keygreen, only marks, mark=o, mark size=1.1pt] coordinates {(0.000,1.188) (0.857,2.513) (1.356,3.233) (1.993,3.645) (2.227,4.021) (2.495,4.242) (2.647,4.344) (2.569,4.373) (2.564,4.355) (2.633,4.290) (2.823,4.546) (2.840,4.401) (2.952,4.566) (2.925,4.584) (2.984,4.568) (2.860,4.535) (3.055,4.623) (3.043,4.622) (3.033,4.634) (2.963,4.624) (2.961,4.598) (2.958,4.604) (2.751,4.557) (2.934,4.606) (2.881,4.599) (3.049,4.628) (3.106,4.608) (2.997,4.633) (2.968,4.565) (3.137,4.659) (3.005,4.639) (3.056,4.650) (3.088,4.670) (2.999,4.656) (3.011,4.639)};
\addplot[color=warnred, only marks, mark=o, mark size=1.1pt] coordinates {(0.000,1.188) (0.855,2.541) (1.397,3.264) (1.992,3.639) (2.220,3.818) (2.730,4.212) (3.141,4.407) (3.486,4.510) (3.768,4.614) (3.985,4.604) (4.307,4.704) (4.389,4.705) (4.419,4.735) (4.486,4.773) (4.653,4.781) (4.693,4.796) (4.568,4.779) (4.569,4.809) (4.770,4.761) (4.937,4.674) (4.716,4.777) (4.725,4.802) (4.694,4.801) (4.668,4.766) (4.714,4.788) (4.690,4.795) (4.770,4.812) (4.672,4.529) (4.765,4.800) (4.829,4.810) (4.833,4.825) (4.819,4.821) (4.820,4.816) (4.801,4.829) (4.832,4.821)};
\addplot[color=sysblue, only marks, mark=*, mark size=1.1pt] coordinates {(0.000,-0.048) (0.817,0.973) (1.229,1.973) (1.477,2.355) (1.915,3.122) (1.815,3.278) (1.855,2.880) (1.934,3.389) (1.840,3.061) (1.942,3.133) (1.873,2.858) (1.860,3.209) (1.923,3.350) (2.020,3.798) (1.964,3.291) (1.829,3.192) (1.856,2.880) (1.993,3.327) (1.893,3.542) (1.918,3.192) (1.920,3.327) (1.817,3.178) (1.806,3.145) (1.944,3.383) (1.872,3.620) (1.952,3.123) (1.914,3.164) (1.807,2.950) (1.892,3.228) (1.965,3.211) (1.883,3.017) (2.089,3.567) (1.881,3.281) (1.881,3.233) (2.018,3.730)};
\addplot[color=keygreen, only marks, mark=*, mark size=1.1pt] coordinates {(0.000,-0.048) (0.857,1.055) (1.356,2.153) (1.993,2.488) (2.227,3.261) (2.495,3.827) (2.647,3.488) (2.569,3.952) (2.564,3.403) (2.633,3.923) (2.823,3.866) (2.840,3.778) (2.952,3.908) (2.925,4.030) (2.984,3.808) (2.860,4.159) (3.055,4.358) (3.043,4.189) (3.033,4.200) (2.963,4.531) (2.961,4.275) (2.958,4.414) (2.751,4.408) (2.934,4.502) (2.881,4.675) (3.049,4.442) (3.106,4.588) (2.997,4.591) (2.968,4.586) (3.137,4.242) (3.005,4.358) (3.056,4.405) (3.088,4.319) (2.999,4.638) (3.011,4.487)};
\addplot[color=warnred, only marks, mark=*, mark size=1.1pt] coordinates {(0.000,-0.048) (0.855,1.044) (1.397,2.095) (1.992,2.845) (2.220,2.867) (2.730,3.614) (3.141,3.498) (3.486,3.280) (3.768,2.891) (3.985,2.392) (4.307,2.345) (4.389,1.803) (4.419,2.077) (4.486,2.055) (4.653,1.767) (4.693,1.859) (4.568,2.152) (4.569,2.247) (4.770,2.159) (4.937,1.678) (4.716,2.367) (4.725,2.513) (4.694,2.598) (4.668,2.780) (4.714,2.720) (4.690,2.847) (4.770,2.806) (4.672,2.517) (4.765,2.809) (4.829,2.763) (4.833,2.825) (4.819,2.895) (4.820,2.842) (4.801,2.872) (4.832,2.817)};
\legend{{proxy, $\beta=0.2$},{proxy, $\beta=0.05$},{proxy, $\beta=0$},{true, $\beta=0.2$},{true, $\beta=0.05$},{true, $\beta=0$}}
\end{axis}
\end{tikzpicture}
\caption{The same runs plotted against $\sqrt{\mathrm{KL}}$, as in Gao et al.\ (2023). Hollow markers are
the proxy (reward-model) score and filled markers the true score. The proxy rises monotonically with KL;
the true score rises, peaks around 3 nats$^{1/2}$ (about 9 nats), and falls beyond it. Up to that
point the KL budget buys real quality; past it, it buys exploits.}
\label{fig:rlhf:sqrtkl}
\end{figure}


Four observations, each of which recurs at scale:

\begin{enumerate}
  \item \textbf{PPO works.} In every run the proxy reward rises from $1.2$ to $4$--$4.8$ within
        about 20 iterations, and responses grow from about 5 to about 7 words. For reference, the best
achievable true score is $7.5$ (six distinct positive words including the topic word); no run gets
close, which leaves room for the lab's stretch goals.
  \item \textbf{The KL penalty sets the operating point.} At convergence the policy spends roughly 3.5,
        9 and 23 nats per response for $\beta=0.2$, $0.05$ and $0$. The best true score comes from the
        middle setting.
  \item \textbf{Without the anchor, the policy over-optimises.} With $\beta=0$ the true score peaks at $3.6$
        (iteration 11) and then decays to about $2.8$, while the proxy is the highest of the three runs.
        \Cref{fig:rlhf:sqrtkl} shows the shape reported by Gao et al.\ (2023) for real reward models.
  \item \textbf{The critic needs time.} Explained variance starts negative ($-0.2$) and climbs to
        $0.6$--$0.9$ over the first 20--40 iterations; early updates rely mostly on the baseline effect.
\end{enumerate}

The final samples make the mechanism visible. For topic \texttt{t2}, the $\beta=0.05$ policy writes
\texttt{p4 n6 p1 p2 p5 p3}: six distinct words, mostly positive, including the topic word \texttt{p2}.
The $\beta=0$ policy has collapsed to one response for every prompt, \texttt{p5 n6 p1 p0 p3 p5 p5}.
It repeats \texttt{p5} because the reward model never learned that repetition is bad, and it has
dropped the topic word because a generic high-reward string is easier to find than a prompt-specific
one. This loss of diversity is itself a symptom worth monitoring (entropy and distinct-$n$ metrics).

\begin{systemslens}[Where the time went on a CPU]
In these runs the learn phase took $87\%$ of each iteration, scoring $8\%$ and generation $5\%$, in
line with the FLOP model of \cref{sec:rlhf:system}, which puts learning at about $83\%$ of FLOPs for
$E=4$ epochs. A 45k-parameter model decodes at full speed on a CPU, because its weights fit in cache. Scale the
policy to billions of parameters on a GPU and generation becomes memory-bound and serial, and the shares invert:
generation is commonly reported as the majority of wall-clock time in PPO and GRPO training. The lab's
measurement task asks you to reproduce the CPU breakdown and to predict the GPU one with the roofline
model of \cref{ch:inference}.
\end{systemslens}

\section{Reward hacking and over-optimisation}\label{sec:rlhf:overopt}

A reward model is a proxy, and optimisation pressure finds the places where the proxy and the
target disagree, a direct instance of Goodhart's law. Gao et al.\ (2023) measured this systematically
by training policies against a proxy reward model and scoring them with a larger ``gold'' reward model. As the KL
from the initial policy grows, the gold score first rises with the proxy, then peaks and falls,
while the proxy keeps rising. The gap widens faster for smaller reward models and less data.
Measuring progress against $\sqrt{\KL}$ rather than steps makes these curves comparable across
optimisers (PPO and best-of-$n$ sampling show the same shape).

Common symptoms in real systems:
\begin{itemize}
  \item \textbf{Length and formatting inflation.} Annotators like thorough answers; the policy learns
        that longer is always better. Responses fill with bullet points, headers and restatements.
  \item \textbf{Sycophancy.} Agreeing with the user's stated view is rewarded more often than
        correcting it.
  \item \textbf{Degenerate exploits.} Repeated phrases, unusual tokens or formatting quirks that
        happen to score highly, like the repetition exploit in our toy task.
  \item \textbf{Calibration loss.} An RLHF-trained model's token probabilities no longer reflect its
        accuracy, because PPO optimised them for reward, not for likelihood. A base model can be well
        calibrated on multiple-choice questions while its RLHF version is overconfident.
\end{itemize}

Mitigations, in rough order of how often they are used: a KL penalty with a target budget; early
stopping by KL or by a held-out gold evaluation; collecting fresh comparisons on the \emph{current}
policy's samples and retraining the reward model (iterated RLHF, as in Llama~2's five rounds); reward-model ensembles
whose disagreement down-weights uncertain rewards; explicit length penalties or length-controlled
evaluation; and mixing in verifiable rewards (\cref{ch:grpo}) wherever a program can check the
answer.

\begin{keyidea}[Best-of-$n$ as a baseline]
Before running PPO, sample $n$ responses from the SFT model and keep the one the reward model
likes best. Its KL from the SFT policy is known in closed form, $\log n - \frac{n-1}{n}$
(an upper bound, exact for continuous rewards), and it needs no training. If PPO cannot beat best-of-$n$ at the
same KL, something is wrong. Rejection-sampling fine-tuning, which fine-tunes on the best-of-$n$
samples (\cref{ch:pg}), is the training-time version and was used alongside PPO in Llama~2.
\end{keyidea}

\section{The four-model system}\label{sec:rlhf:system}

\Cref{sec:grpo:rlhf} costed the four models at a glance. Here we plan them properly, because
placement determines whether PPO-RLHF fits at all, and how fast it runs.

\begin{table}[ht]
\centering\small
\begin{tabular}{@{}lrrrr@{}}
\toprule
\textbf{Policy size} & \textbf{policy (trained)} & \textbf{value (trained)} & \textbf{ref + RM (bf16)} & \textbf{total}\\
\midrule
1B  & 16\,GB   & 16\,GB   & 4\,GB   & 36\,GB\\
8B  & 128\,GB  & 128\,GB  & 32\,GB  & 288\,GB\\
70B & 1{,}120\,GB & 1{,}120\,GB & 280\,GB & 2{,}520\,GB\\
\bottomrule
\end{tabular}
\caption{Static memory of PPO-RLHF when all four models are the policy's size (16 bytes per trained
parameter for mixed-precision AdamW, 2 bytes per frozen parameter). This excludes activations, the KV
cache for generation and communication buffers. Computed with \code{llmsys.rlhf.plan.state\_bytes}.}
\label{tab:rlhf:memory}
\end{table}

\begin{napkinmath}[How many GPUs for an 8B PPO-RLHF run?]
From \cref{tab:rlhf:memory}, 288\,GB of static state. With full sharding (ZeRO-3/FSDP,
\cref{ch:sharding}) over 80\,GB GPUs at 70\% usable memory (56\,GB each), the minimum is
$\lceil 288/56\rceil = 6$ GPUs, and in practice 8, which leaves headroom for activations and the
rollout KV cache. That cache is substantial: for Qwen3-8B (36 layers, 8 KV heads of dimension 128, bf16) it is
$2\cdot36\cdot8\cdot128\cdot2 = 147{,}456$ bytes per token. 512 responses of 1,024 tokens
therefore need $512\cdot1024\cdot147{,}456\approx 77$\,GB of KV cache, about one more GPU's worth during
the rollout phase alone. \code{plan.min\_gpus} and \code{plan.kv\_cache\_bytes} reproduce these numbers.
\end{napkinmath}

\begin{napkinmath}[Where the FLOPs go in one iteration]
Take $B=512$ responses of $T=1{,}024$ tokens (prompt plus response), $N$ parameters in each of the
four models, and $E=2$ PPO epochs. With the $2N$/$6N$ rule (\cref{ch:pretrain}):
generation $\approx 2NBT$; scoring $\approx 4\cdot 2NBT = 8NBT$; learning $\approx E\cdot 6\cdot 2N\cdot BT = 24NBT$.
Learning dominates the FLOPs (about 70\%), yet in wall-clock terms generation is often the largest phase,
because it runs at a small fraction of peak FLOP/s (\cref{ch:inference}). This is why every serious RLHF
system replaces the trainer's \code{generate()} with an inference engine (\cref{ch:serving}).
\end{napkinmath}

\begin{figure}[ht]
\centering
\resizebox{\linewidth}{!}{%
\begin{tikzpicture}[
  gpu/.style={draw, rounded corners=2pt, minimum width=21mm, minimum height=24mm, fill=black!3},
  m/.style={draw, rounded corners=1pt, minimum width=18mm, minimum height=4.5mm, font=\tiny},
  >={Stealth[length=2mm]}]
\node[font=\small\bfseries] at (2.4,2.0) {(a) colocated, time-sliced};
\foreach \g in {0,1} {
  \node[gpu] (A\g) at (1.2+2.4*\g,0) {};
  \node[font=\tiny] at (1.2+2.4*\g,-1.45) {GPU \g};
  \node[m, fill=fcompute!35] at (1.2+2.4*\g,0.8) {policy shard};
  \node[m, fill=fcompute!20] at (1.2+2.4*\g,0.25) {value shard};
  \node[m, fill=black!12] at (1.2+2.4*\g,-0.3) {ref shard};
  \node[m, fill=napkin!30] at (1.2+2.4*\g,-0.85) {RM shard};
}
\node[font=\small\bfseries] at (10.2,2.0) {(b) split placement};
\foreach \g/\lbl/\c in {0/{policy}/fcompute!35, 1/{policy}/fcompute!35, 2/{value}/fcompute!20, 3/{ref + RM}/napkin!30} {
  \node[gpu] (B\g) at (7.0+2.3*\g,0) {};
  \node[m, fill=\c, minimum height=14mm, text width=16mm, align=center] at (7.0+2.3*\g,0) {\lbl};
  \node[font=\tiny] at (7.0+2.3*\g,-1.45) {GPU \the\numexpr\g+2\relax};
}
\node[draw, rounded corners, fill=fmemory!25, font=\scriptsize, align=center] (eng) at (7.0,-2.4)
  {rollout engine\\(vLLM / SGLang)};
\draw[->, dashed] (B0.south) -- node[right, font=\tiny]{weight sync} (eng.north);
\node[font=\scriptsize, text width=48mm, align=left] at (2.4,-2.6) {All models sharded on all GPUs;
phases run one after another; idle models can be offloaded to CPU memory between phases.};
\node[font=\scriptsize, text width=50mm, align=left] at (12.0,-2.6) {Each model on its own GPU group;
scoring and learning can be pipelined, but some groups idle while others work.};
\end{tikzpicture}}
\caption{Two ways to place the four models. Colocation (a) maximises memory per model and avoids
idle hardware; it needs fast switching between phases (offload, resharding). Split placement (b),
as in Ray-based frameworks such as OpenRLHF, scales each role independently; the
trade-offs are analysed in \cref{sec:rlsys:placement}.}
\label{fig:rlhf:placement}
\end{figure}

\begin{systemslens}[Tricks that shrink the four-model tax]
\begin{itemize}
  \item \textbf{Reference via LoRA.} If the policy is trained with LoRA adapters (\cref{sec:instruct:lora}),
        the reference model is the same network with the adapters disabled, so no second copy is needed.
  \item \textbf{Precompute the reference.} $\log\piref$ depends only on the sampled tokens, so it is computed
        once per rollout (phase 2) and never again; the reference can live on CPU or another node.
  \item \textbf{Shared trunk critic.} One trunk with two heads halves the trained state, at a
        stability cost (\cref{sec:rlhf:value}).
  \item \textbf{Smaller reward and value models.} InstructGPT's 6B reward model for a 175B policy is the
        canonical example.
  \item \textbf{Drop the critic.} GRPO and RLOO replace $V_\psi$ with group statistics
        (\cref{ch:grpo}), which removes 16 bytes per parameter of trained state.
  \item \textbf{Drop the reward model.} Verifiable rewards (RLVR) replace $r_\phi$ with a program.
\end{itemize}
\end{systemslens}

\begin{lstlisting}[caption={Planning helpers (\ident{llmsys/rlhf/plan.py}).},label={lst:rlhf:plan}]
TRAINED_BYTES, FROZEN_BYTES = 16, 2

def state_bytes(m, offload_frozen=False, zero_shards=1):
    frozen = 0.0 if offload_frozen else FROZEN_BYTES
    out = {"policy": m.policy * TRAINED_BYTES / zero_shards,
           "value": m.value * TRAINED_BYTES / zero_shards,
           "reward": m.reward * frozen / zero_shards,
           "reference": m.reference * frozen / zero_shards}
    out["total"] = sum(out.values())
    return out

>>> state_bytes(RLHFModels(8e9, 8e9, 8e9, 8e9))["total"] / 1e9
288.0
\end{lstlisting}

\section{RL from AI feedback}\label{sec:rlhf:rlaif}

Human comparisons are slow, expensive and inconsistent. \emph{RL from AI feedback} (RLAIF) replaces
some or all of the annotators with an LLM judge. Two variants are common:

\begin{enumerate}
  \item \textbf{AI-labelled preferences, then a reward model.} A strong model compares two responses
        under written instructions and produces the pairwise labels; the rest of the pipeline is unchanged.
        Constitutional AI (Bai et al., 2022) wrote those instructions as a short list of principles
        (the ``constitution''). It also added a supervised stage in which the model critiques and revises its own
        responses according to sampled principles.
  \item \textbf{Direct scoring.} Skip the reward model and ask the judge for a score at PPO time. This is
        simpler but expensive: a judge call per rollout, in the training loop's critical path.
\end{enumerate}

LLM judges agree with humans about as often as humans agree with each other on many
general-domain tasks, but they bring systematic biases: preferring the first-listed response
(position bias), preferring their own model family, and preferring longer answers. The mitigations are
those of \cref{sec:eval:judge}: evaluate both orders and average, use rubrics, calibrate against a
human-labelled subset, and track agreement over time. The systems benefit is large: an AI labeller
is a batch inference job that can produce millions of comparisons on demand, so preference data can
be refreshed every RLHF round on the current policy's samples.

\section{Agentic RL with milestone rewards}\label{sec:rlhf:agentic}

When the policy is an agent (calling tools, browsing, editing files over many turns), the episode
is a multi-turn trajectory and a single terminal reward is extremely sparse: did a 40-step
coding session end with passing tests? \emph{Milestone rewards} add intermediate credit for verifiable
sub-goals, for example: the repository was cloned, the failing test was located, a patch compiles, the
target test passes, no other test broke. Each milestone is checked by a program, and its reward is
attached to the token that completed the corresponding action, which turns the shaped-reward
machinery of \cref{eq:rlhf:shaped} into a natural fit.

\begin{itemize}
  \item \textbf{Mask what the policy did not write.} Tool outputs and environment observations are part of
        the context but not actions; their tokens get mask 0 in every loss (\cref{sec:rlsys:agentic}).
  \item \textbf{Values per turn or per token.} A critic can still value every token, but GAE across
        long tool outputs is wasteful; some systems compute advantages per turn.
  \item \textbf{Guard against milestone farming.} Pay each milestone at most once per episode, and keep the
        final-success reward dominant, or the agent will learn to collect easy milestones and stop.
\end{itemize}

\begin{lstlisting}[caption={Milestone rewards attached to the tokens that completed them.},label={lst:rlhf:milestones}]
MILESTONES = {"tests_located": 0.1, "patch_compiles": 0.2, "target_passes": 0.5,
              "suite_green": 1.0}

def milestone_rewards(events, n_tokens):
    """events: (milestone, token_index) pairs from the environment checker."""
    r, seen = torch.zeros(n_tokens), set()
    for name, t in events:
        if name in MILESTONES and name not in seen:   # pay each one once
            r[t] += MILESTONES[name]
            seen.add(name)
    return r          # add to the KL-shaped per-token reward before GAE
\end{lstlisting}

\section{Case studies}\label{sec:rlhf:cases}

\paragraph{InstructGPT (2022).} SFT on about 13k demonstrations; a 6B reward model on about 33k prompts ranked with $K=4$--$9$
responses; PPO with a per-token KL penalty and PPO-ptx. Labellers preferred the 1.3B InstructGPT model
to the 175B GPT-3, which showed that RLHF can be worth more than a 100$\times$ increase in parameters
for following instructions.

\paragraph{Llama~2 (2023).} Two separate reward models, one for helpfulness and one for safety, with
a margin term from the annotators' preference strength. Five rounds of iterated RLHF, each collecting new comparisons on
the latest model; the early rounds used rejection-sampling fine-tuning (best-of-$n$ under the reward
model), later combined with PPO. The two reward scores were combined with a rule that prioritises
safety on prompts likely to be unsafe.

\paragraph{After Llama~2.} Many recent open models replace PPO with offline or online DPO-style
objectives (\cref{ch:dpo}) for general chat, and use RL with verifiable rewards (\cref{ch:grpo}) for
maths and code. Meta's Llama~4 announcement, for example, described a lightweight SFT, then online RL, then lightweight DPO
sequence. Treat such summaries as descriptions of a recipe's shape rather than of its details.

\begin{table}[ht]
\centering\small
\begin{tabularx}{\linewidth}{@{}lXXXX@{}}
\toprule
& \textbf{PPO-RLHF} & \textbf{GRPO / RLVR} & \textbf{DPO} & \textbf{Best-of-$n$ / RFT}\\
\midrule
Data & prompts + RM & prompts + verifier & fixed preference pairs & prompts + RM\\
Online sampling & yes & yes (groups) & no (offline) & yes (at inference or once)\\
Models in memory & 4 (2 trained) & 1--2 (1 trained) & 2 (1 trained) & 1 + RM\\
Baseline & learned critic & group mean & implicit & none\\
Credit assignment & per token (GAE) & per sequence & per sequence & per sequence\\
Main risk & instability, cost & sparse reward & off-distribution data & inference cost\\
Chapter & this one & \ref{ch:grpo}--\ref{ch:grpo2} & \ref{ch:dpo} & \ref{ch:pg}, \ref{ch:scaling}\\
\bottomrule
\end{tabularx}
\caption{The main post-training optimisers side by side.}
\label{tab:rlhf:compare}
\end{table}

\section{Summary}\label{sec:rlhf:summary}
\begin{itemize}
  \item RLHF optimises a learned reward under a KL anchor to the SFT model,
        \cref{eq:rlhf:objective}. The anchor is what keeps a proxy reward useful.
  \item A reward model is the policy architecture with a scalar head, trained with the Bradley--Terry loss
        on comparisons. Read the score at each row's last real token, group $K$-way rankings by prompt,
        and expect held-out accuracy near annotator agreement.
  \item PPO-RLHF turns one terminal score into per-token learning signals: KL-shaped rewards,
        a learned critic, and GAE's backward scan, $\hat A_t=\delta_t+\gamma\lambda\hat A_{t+1}$.
  \item The update is the clipped surrogate for the actor plus a (clipped) regression for the critic,
        for a few epochs over each rollout batch, with rollout-time log-probabilities held fixed.
  \item Over-optimisation is the default outcome of optimising a proxy hard enough. Watch true-quality
        evaluations against KL, not reward against steps.
  \item The system cost is four models, two of them trained, plus autoregressive generation. Memory planning,
        placement, and offloading the reference and reward models decide feasibility; generation usually decides speed.
\end{itemize}

\section{Exercises}\label{sec:rlhf:exercises}

\begin{exercise}[BT from Gumbel noise]
Show that if $u_w = r_w + \epsilon_w$ and $u_l = r_l + \epsilon_l$ with independent standard Gumbel noise,
then $\Pr[u_w > u_l] = \sigma(r_w - r_l)$. What changes if the noise is Gaussian?
\end{exercise}

\begin{exercise}[The gradient of the BT loss]
Compute $\partial\mathcal{L}_{\text{RM}}/\partial r_w$ for a single pair and interpret it. Which pairs
dominate the gradient late in training? How does a margin $m$ change this?
\end{exercise}

\begin{exercise}[Grouped rankings]
For $K=6$ ranked responses, how many forward passes does the grouped ranking loss need per prompt, and how
many would independent pairs need? Over one epoch, how many gradient steps does each response contribute
to in each case?
\end{exercise}

\begin{exercise}[GAE limits]
Using \cref{eq:rlhf:gae}, show that $\lambda=1,\gamma=1$ gives $\hat A_t=\sum_{k\ge t}r_k-V(s_t)$ and that
$\lambda=0$ gives $\hat A_t=\delta_t$. Recompute the $\lambda=0.5$ row of \cref{tab:rlhf:gae} by hand,
then check it with \code{gae()}.
\end{exercise}

\begin{exercise}[KL estimators]
For $\piold$ and $\piref$ both Bernoulli, compute the exact $\KL(\piold\|\piref)$ and the expectations and
variances of the $k_1$, $k_2$ and $k_3$ estimators under samples from $\piold$, for
$p_{\text{old}}=0.6$ and $p_{\text{ref}}\in\{0.5,0.2\}$. Which estimators are unbiased?
\end{exercise}

\begin{exercise}[Off-by-one critic]
Modify the trainer to read values at positions $P,\dots,P+T-1$ instead of $P-1,\dots,P+T-2$. Predict
what happens to explained variance and to the reward curve, then run it and compare.
\end{exercise}

\begin{exercise}[Memory planning]
Using \code{plan.state\_bytes}, compute the static memory for a 32B policy with a 7B reward model, a
critic initialised from that reward model, and a LoRA reference. How many 80\,GB GPUs are needed with full
sharding at 70\% usable memory?
\end{exercise}

\begin{exercise}[Best-of-$n$ KL]
Show that the KL between the best-of-$n$ distribution and the base policy is at most $\log n-\frac{n-1}{n}$.
For which $n$ does best-of-$n$ spend the same KL as the PPO run in the lab?
\end{exercise}

\begin{exercise}[Adaptive KL]
Run the lab with \code{target\_kl} set to 1, 3 and 10 nats per response. Plot true score against final
KL. Where does over-optimisation set in?
\end{exercise}

\begin{exercise}[Design: RLAIF judge in the loop]
A judge model of 70B parameters scores every rollout of an 8B policy directly (no reward model).
Estimate the judge's share of each iteration's FLOPs and wall time for $B=512$ responses of 512 tokens with a
256-token rubric prompt. Propose two system changes that remove it from the critical path.
\end{exercise}

\section{Lab 35: PPO-RLHF from scratch}\label{sec:rlhf:lab}

\begin{labbox}{35}{PPO-RLHF from scratch}
\labdir{lab35_rlhf_ppo}
\labsection{Goals}
Implement the mathematical core of RLHF as small functions, plug them into the provided four-model trainer,
and train a policy against a learned reward model on a laptop CPU. Then watch PPO raise the proxy
reward, and push it hard enough to over-optimise.

\labsection{Setup}
\begin{lstlisting}[style=shell]
cd code
LAB_IMPL=starter  python -m pytest labs/lab35_rlhf_ppo -q   # your code
LAB_IMPL=solution python -m pytest labs/lab35_rlhf_ppo -q   # reference
python labs/lab35_rlhf_ppo/run.py all --iters 70   # pipeline + plots
\end{lstlisting}
The tests use the synthetic preference task of \cref{sec:rlhf:results} and tiny models (untied
embeddings), so the full suite runs in under a minute on a CPU.

\labsection{Tasks}
\begin{enumerate}
  \item \code{masked\_mean}, \code{whiten}: statistics over real tokens only.
  \item \code{bt\_loss}, \code{rm\_accuracy}: the Bradley--Terry loss, with an optional margin.
  \item \code{ranking\_loss}: the grouped $K$-way loss; it must equal \code{bt\_loss} when $K=2$.
  \item \code{kl\_shaped\_rewards}: per-token KL tolls plus the score on each row's last real token.
  \item \code{gae}: the backward scan with masking; tested against a brute-force sum of discounted TD residuals.
  \item \code{ppo\_policy\_loss}: the clipped surrogate with \code{clipfrac} and \code{approx\_kl} statistics.
  \item \code{value\_loss}: the (optionally clipped) critic loss.
  \item \code{explained\_variance} and \code{AdaptiveKLController.update}.
  \item Run \code{train\_pipeline()} with your functions: SFT, preference sampling and labelling, reward-model
        training, then PPO. The smoke test checks that the reward model beats chance on held-out pairs and
        that the reward rises.
\end{enumerate}

\labsection{Measurements}
\begin{itemize}
  \item Plot RM score, true score, KL, explained variance and response length against iteration for
        $\beta\in\{0.2, 0.05, 0\}$. Find the iteration where true score peaks for $\beta=0$.
  \item Compare $\lambda\in\{0, 0.95, 1\}$: the speed of reward improvement and the critic's explained variance.
  \item Time each phase (rollout, score, learn) and compare the shares with the FLOP model of
        \cref{sec:rlhf:system}. On the CPU, which phase dominates, and why might it differ on a GPU?
\end{itemize}

\labsection{Deliverables}
Passing tests; a figure like \cref{fig:rlhf:curves} for your own runs; a paragraph explaining the
over-optimisation you observe in terms of what the reward model never saw.

\labsection{Stretch goals}
Train a reward-model ensemble and penalise disagreement; add PPO-ptx with the SFT data; replace the reward
model with the hidden score (an ``oracle'' run) and compare learning speed; implement per-sequence loss
aggregation and compare length dynamics.
\end{labbox}
